\documentclass[11pt,a4paper]{article}
% Preâmbulo comum. Não duplicar pacotes nos módulos.
\usepackage[utf8]{inputenc}
\usepackage[T1]{fontenc}
\usepackage[brazil]{babel}
\usepackage{amsmath,amssymb,amsthm,mathtools}
\usepackage{geometry}
\geometry{a4paper,margin=2.3cm}
\usepackage{microtype}
\usepackage{enumitem}
\usepackage{graphicx}
\usepackage{xcolor}
\usepackage{booktabs}
\usepackage{hyperref}
\hypersetup{colorlinks=true,linkcolor=blue!40!black,urlcolor=blue!50!black}

\setcounter{tocdepth}{2}
\numberwithin{equation}{section}

\theoremstyle{plain}
\newtheorem{teorema}{Teorema}[section]
\newtheorem{proposicao}[teorema]{Proposição}
\newtheorem{lema}[teorema]{Lema}
\newtheorem{corolario}[teorema]{Corolário}

\theoremstyle{definition}
\newtheorem{definicao}[teorema]{Definição}
\newtheorem{exemplo}[teorema]{Exemplo}
\newtheorem{exercicio}{Exercício}[section]

\theoremstyle{remark}
\newtheorem{observacao}[teorema]{Observação}
\newtheorem{contraexemplo}[teorema]{Contraexemplo}
\newtheorem{intuicao}[teorema]{Intuição}

\newcommand{\R}{\mathbb{R}}
\newcommand{\N}{\mathbb{N}}
\newcommand{\Z}{\mathbb{Z}}
\newcommand{\Q}{\mathbb{Q}}
\DeclareRobustCommand{\eps}{\varepsilon}
\DeclareRobustCommand{\vet}[1]{\vec{#1}}
\DeclareRobustCommand{\dif}{\mathrm{d}}
\DeclareRobustCommand{\abs}[1]{\left\lvert #1 \right\rvert}
\DeclareRobustCommand{\norm}[1]{\left\lVert #1 \right\rVert}
\DeclareMathOperator{\grad}{grad}
\DeclareMathOperator{\Div}{div}
\DeclareMathOperator{\rot}{rot}

\newcommand{\base}[1]{\par\medskip\noindent\textbf{Base para relembrar.} #1\par\medskip}
\newcommand{\cuidado}[1]{\par\medskip\noindent\textbf{Quando a hipótese falha.} #1\par\medskip}

\newcommand{\figura}[2]{%
  \begin{figure}[htbp]
  \centering
  \IfFileExists{#1}{%
    \includegraphics[width=0.84\textwidth]{#1}%
  }{%
    \fbox{\parbox{0.8\textwidth}{\centering Figura ainda não gerada:\\ \texttt{#1}}}%
  }
  \caption{#2}
  \end{figure}%
}

\newenvironment{esqueleto}{%
  \par\medskip\noindent\textbf{Esqueleto da demonstração.}\par\nopagebreak
  \begin{enumerate}[leftmargin=1.4em]
}{\end{enumerate}\medskip}

\title{Gabarito: Green, Stokes e Gauss}
\author{}
\date{}
\begin{document}
\maketitle
\tableofcontents

\section{Exercícios básicos}

\begin{exercicio}[{[Básico --- Green num retângulo]}]
Seja $D=[0,1]\times[0,2]$, com fronteira anti-horária, $P=-y$ e $Q=x$. Calcule os dois lados da fórmula de Green.
\end{exercicio}

\noindent\textbf{Solução.}
$\partial_x Q-\partial_y P=1-(-1)=2$. A integral dupla é
\[
\iint_D 2\,\dif A=2\cdot 1\cdot 2=4.
\]
Os quatro lados:
\begin{itemize}
\item Base, $y=0$, $x$ de $0$ a $1$: $P=0$ e $\dif y=0$, contribuição $0$.
\item Direita, $x=1$, $y$ de $0$ a $2$: $\dif x=0$ e $Q=1$, contribuição $\int_0^2 1\,\dif y=2$.
\item Topo, $y=2$, $x$ de $1$ a $0$: $P=-2$ e $\dif y=0$, contribuição $\int_1^0 -2\,\dif x=2$.
\item Esquerda, $x=0$, $y$ de $2$ a $0$: $\dif x=0$ e $Q=0$, contribuição $0$.
\end{itemize}
A circulação vale $4$, igual à integral dupla.

\begin{exercicio}[{[Básico --- divergência e rotacional]}]
Para $\vet{F}=(xy,\, yz,\, zx)$, calcule $\nabla\cdot\vet{F}$ e $\nabla\times\vet{F}$.
\end{exercicio}

\noindent\textbf{Solução.}
Com $P=xy$, $Q=yz$ e $R=zx$,
\[
\nabla\cdot\vet{F}=\partial_x(xy)+\partial_y(yz)+\partial_z(zx)=y+z+x.
\]
O rotacional é
\begin{align*}
\partial_y R-\partial_z Q&=\partial_y(zx)-\partial_z(yz)=0-y=-y,\\
\partial_z P-\partial_x R&=\partial_z(xy)-\partial_x(zx)=0-z=-z,\\
\partial_x Q-\partial_y P&=\partial_x(yz)-\partial_y(xy)=0-x=-x.
\end{align*}
Logo $\nabla\times\vet{F}=(-y,-z,-x)$.

\begin{exercicio}[{[Básico --- área do círculo]}]
Prove, por Green, que o disco de raio $R>0$ tem área $\pi R^2$.
\end{exercicio}

\noindent\textbf{Solução.}
A fronteira $x=R\cos t$, $y=R\sin t$, $t\in[0,2\pi]$, é anti-horária. Então $\dif x=-R\sin t\,\dif t$ e $\dif y=R\cos t\,\dif t$, e
\begin{align*}
-y\,\dif x+x\,\dif y
&=-(R\sin t)(-R\sin t\,\dif t)+(R\cos t)(R\cos t\,\dif t)\\
&=R^2(\sin^2 t+\cos^2 t)\,\dif t
=R^2\,\dif t.
\end{align*}
A terceira fórmula da área dá
\[
\frac12\int_0^{2\pi} R^2\,\dif t=\pi R^2.
\]
$P=-y/2$ e $Q=x/2$ são $C^1$ e o disco é região dos tipos~I e~II, então Green se aplica. De fato $\partial_x Q-\partial_y P=1/2-(-1/2)=1$.

\begin{exercicio}[{[Básico --- área e fluxo de um gráfico]}]
Seja $\vet{r}(u,v)=(u,v,u+v)$ em $[0,1]^2$ e $\vet{F}=(0,0,1)$. Calcule o produto fundamental, a área e o fluxo.
\end{exercicio}

\noindent\textbf{Solução.}
$\vet{r}_u=(1,0,1)$ e $\vet{r}_v=(0,1,1)$. O produto vetorial é
\begin{align*}
\bigl(0\cdot 1-1\cdot 1,\;
1\cdot 0-1\cdot 1,\;
1\cdot 1-0\cdot 0\bigr)
=(-1,-1,1).
\end{align*}
A norma vale $\sqrt{3}$, e a área é $\int_0^1\int_0^1\sqrt{3}\,\dif u\,\dif v=\sqrt{3}$. O campo $\vet{F}=(0,0,1)$ tem produto escalar $1$ com o produto fundamental, então o fluxo é
\[
\int_0^1\int_0^1 1\,\dif u\,\dif v=1.
\]
A componente $z$ da normal é positiva.

\section{Exercícios intermediários}

\begin{exercicio}[{[Intermediário --- circulação pelo Green]}]
Calcule $\oint_C -y^2\,\dif x+x\,\dif y$ na fronteira anti-horária do quadrado $[0,1]^2$.
\end{exercicio}

\noindent\textbf{Solução.}
$P=-y^2$ e $Q=x$ são de classe $C^1$ em $\R^2$, e o quadrado é dos tipos~I e~II. Green se aplica:
\[
\partial_x Q-\partial_y P=1-(-2y)=1+2y.
\]
A integral dupla é
\[
\int_0^1\int_0^1(1+2y)\,\dif x\,\dif y
=\int_0^1(1+2y)\,\dif y
=\bigl[y+y^2\bigr]_0^1=2.
\]
Conferência pelos lados. Na base $y=0$, $P=0$. Na direita $x=1$, $y$ de $0$ a $1$, $Q=1$ e $\int_0^1 1\,\dif y=1$. No topo $y=1$, $x$ de $1$ a $0$, $P=-1$ e $\int_1^0 -1\,\dif x=1$. Na esquerda $x=0$, $Q=0$. A soma é $2$.

\begin{exercicio}[{[Intermediário --- Stokes no disco]}]
Calcule a circulação de $\vet{F}=(z,x,y)$ na circunferência unitária do plano $z=0$, anti-horária, pelo teorema de Stokes, e confira pela linha.
\end{exercicio}

\noindent\textbf{Solução.}
\[
\nabla\times\vet{F}
=\bigl(\partial_y y-\partial_z x,\;\partial_z z-\partial_x y,\;\partial_x x-\partial_y z\bigr)
=(1,1,1).
\]
O disco $x^2+y^2\le 1$, $z=0$, tem normal $\vet{k}$ compatível com o bordo anti-horário. Stokes lê a terceira componente:
\[
\iint_S 1\,\dif A=\pi.
\]
Na linha, $\vet{r}(t)=(\cos t,\sin t,0)$, $\vet{F}(\vet{r})=(0,\cos t,\sin t)$ e $\vet{r}'=(-\sin t,\cos t,0)$. O produto escalar é $\cos^2 t$, e
\[
\int_0^{2\pi}\frac{1+\cos 2t}{2}\,\dif t=\pi.
\]

\begin{exercicio}[{[Intermediário --- Gauss no cubo]}]
Fluxo de $\vet{F}=(x,y,z)$ para fora de $E=[0,1]^3$.
\end{exercicio}

\noindent\textbf{Solução.}
$\nabla\cdot\vet{F}=3$ e o volume do cubo é $1$, então o fluxo externo é $3$.

Nas faces, a normal para fora dá:
\begin{itemize}
\item $x=1$: $\iint_{[0,1]^2} 1\,\dif y\,\dif z=1$;
\item $x=0$: $\iint -x\,\dif y\,\dif z=0$;
\item $y=1$: $\iint 1\,\dif x\,\dif z=1$;
\item $y=0$: contribuição $0$;
\item $z=1$: $\iint 1\,\dif x\,\dif y=1$;
\item $z=0$: contribuição $0$.
\end{itemize}
A soma das seis faces é $3$.

\section{Exercícios avançados}

\begin{exercicio}[{[Avançado --- Green em região tipo II]}]
Prove $\oint Q\,\dif y=\iint\partial_x Q$ numa região do tipo~II, com a fronteira anti-horária.
\end{exercicio}

\noindent\textbf{Solução.}
Seja $D=\{(x,y):c\le y\le d,\; h_1(y)\le x\le h_2(y)\}$, com $h_1\le h_2$ de classe $C^1$ e $Q$ de classe $C^1$ num aberto que contém $D$. A fronteira anti-horária sobe em $x=h_2(y)$, com $y$ de $c$ a $d$, e desce em $x=h_1(y)$, com $y$ de $d$ a $c$. Segmentos horizontais, se existirem nos extremos, têm $\dif y=0$ e não entram nesta integral. Portanto
\begin{align*}
\oint_{\partial D} Q\,\dif y
&=\int_c^d Q\bigl(h_2(y),y\bigr)\,\dif y
+\int_d^c Q\bigl(h_1(y),y\bigr)\,\dif y\\
&=\int_c^d\Bigl(Q\bigl(h_2(y),y\bigr)-Q\bigl(h_1(y),y\bigr)\Bigr)\dif y.
\end{align*}
Para cada $y$, o teorema fundamental do cálculo na variável $x$ produz
\[
Q\bigl(h_2(y),y\bigr)-Q\bigl(h_1(y),y\bigr)
=\int_{h_1(y)}^{h_2(y)}\partial_x Q(x,y)\,\dif x,
\]
porque $\partial_x Q$ é contínua. Logo
\[
\oint_{\partial D} Q\,\dif y
=\int_c^d\int_{h_1(y)}^{h_2(y)}\partial_x Q\,\dif x\,\dif y
=\iint_D\partial_x Q\,\dif A.
\]

\begin{exercicio}[{[Avançado --- Green não vê o furo]}]
Para a forma angular no plano furado, calcule o rotacional escalar, a circulação na circunferência unitária, e diga qual hipótese de Green falha no disco.
\end{exercicio}

\noindent\textbf{Solução.}
Sejam $P=-y(x^2+y^2)^{-1}$ e $Q=x(x^2+y^2)^{-1}$ em $\R^2\setminus\{\vet{0}\}$.
\begin{align*}
\partial_x Q
&=\frac{x^2+y^2-2x^2}{(x^2+y^2)^2}
=\frac{y^2-x^2}{(x^2+y^2)^2},\\
\partial_y P
&=\frac{-(x^2+y^2)+2y^2}{(x^2+y^2)^2}
=\frac{y^2-x^2}{(x^2+y^2)^2}.
\end{align*}
A diferença é zero em todo o domínio.

Na circunferência $\vet{r}(t)=(\cos t,\sin t)$, $t\in[0,2\pi]$, o denominador vale $1$ e $\vet{F}(\vet{r})=(-\sin t,\cos t)=\vet{r}'$. A circulação é $\int_0^{2\pi}1\,\dif t=2\pi$.

Green no disco unitário pediria $P$ e $Q$ de classe $C^1$ num aberto que contém o disco fechado. Todo aberto que contém o disco contém a origem, e na origem a fórmula não está definida: o denominador se anula e o campo não admite extensão contínua, porque $\norm{\vet{F}}=1/\sqrt{x^2+y^2}$ explode. A integral dupla formal de $0$ no disco vale $0$, que não é a circulação $2\pi$. Os números divergem porque a hipótese de regularidade até dentro do furo não se cumpre.

\end{document}
